\subsection{功能性需求分类}

\begin{center}
\renewcommand{\arraystretch}{1.5}
\begin{longtable}{|>{\raggedright\arraybackslash}p{2.8cm}|>{\raggedright\arraybackslash}p{4.2cm}|>{\raggedright\arraybackslash}p{7cm}|}
\hline
\centering 功能类别 & \centering 功能名称、标识符 & \centering 描述 \tabularnewline
\hline
\endfirsthead

\hline
\centering 功能类别 & \centering 功能名称、标识符 & \centering 描述 \tabularnewline
\hline
\endhead

 \legacyold{用户展示模块} & \shortstack[l]{\legacyold{个人信息展示}\\\legacyold{FR-YHZS-0001}} & \legacyold{个人信息展示包括用户昵称、ID、头像、等级等内容。} \tabularnewline
\hline
 \legacyold{用户展示模块} & \shortstack[l]{\legacyold{检测记录展示}\\\legacyold{FR-YHZS-0002}} & \legacyold{按照时间顺序展示用户上传的学术图像、论文与同行评审 Review 检测任务条目，点击可进入检测任务详情页。支持按时间、类型等条件筛选，支持查看与下载检测报告。} \tabularnewline
\hline
 \legacyold{用户展示模块} & \shortstack[l]{\legacyold{审核任务展示}\\\legacyold{FR-YHZS-0003}} & \legacyold{按照时间顺序展示用户提交的人工审核任务条目，支持查看当前审核进度、阶段性结果与最终审核结论。} \tabularnewline
\hline
 \legacyold{用户展示模块} & \shortstack[l]{\legacyold{社区反馈展示}\\\legacyold{FR-YHZS-0004}} & \legacyold{管理员对用户举报或被举报的处理结果、个人评论的新回复将在此处反馈。} \tabularnewline
\hline
用户展示模块 & \shortstack[l]{\iternew{综合鉴伪报告下载与查看}\\\iternew{FR-YHZS-0005}} & \iternew{用户可查看并下载综合鉴伪报告，报告内容包括 AI 贡献占比、图像 AI 可疑位置、使用建议、综合结论等维度。} \tabularnewline
\hline
用户展示模块 & \shortstack[l]{\iternew{人工审核进度与结果查询}\\\iternew{FR-YHZS-0006}} & \iternew{用户可查询人工审核任务的处理进度、参与情况与最终结果，支持对图像、论文和 Review 三类内容进行统一查询。} \tabularnewline
\hline
用户展示模块 & \shortstack[l]{\iternew{检测模型展示与选择}\\\iternew{FR-YHZS-0007}} & \iternew{用户可查看当前检测任务所使用的文本、图像及 Review 检测模型信息，包括模型名称、版本号及检测模式；高级用户可根据需求选择快速检测模式或精准检测模式。} \tabularnewline
\hline

用户上传模块 & \shortstack[l]{文件上传\\FR-YHSC-0001} & 用户可上传待检测的学术图像、全篇论文与同行评审 Review，支持多种格式上传，确保系统兼容性与通用性。 \tabularnewline
\hline
用户上传模块 & \shortstack[l]{AI检测\\FR-YHSC-0002} & 用户上传内容后可获得自动检测结果，系统可对学术图像、全篇论文与同行评审 Review 是否存在造假或 AI 生成倾向进行分析。 \tabularnewline
\hline
用户上传模块 & \shortstack[l]{人工审核请求\\FR-YHSC-0003} & 用户在收到 AI 检测结果后可申请人工审核，任务将进入人工审核流程，等待其他用户参与审查。 \tabularnewline
\hline
用户上传模块 & \shortstack[l]{批量上传\\FR-YHSC-0004} & 用户可批量上传多份学术资源进行统一检测，支持大规模检测场景下的资源提交。 \tabularnewline
\hline
用户上传模块 & \shortstack[l]{高并发批量检测\\FR-YHSC-0005} & 系统提供稳定的多线程上传接口，支持大规模学术资源的并发处理，提高批量检测效率。 \tabularnewline
\hline
用户上传模块 & \shortstack[l]{检测状态实时同步\\FR-YHSC-0006} & 系统支持检测状态的实时同步与更新，用户可及时查看任务处理状态和结果返回情况。 \tabularnewline
\hline

 \legacyold{用户审核模块} & \shortstack[l]{\legacyold{审核任务展示}\\\legacyold{FR-YHSH-0001}} & \legacyold{按照发布时间先后展示用户要求人工审核的任务条目，支持按内容类型、学科等标签分类查看，并在详情页展示已有审核结果。} \tabularnewline
\hline
 \legacyold{用户审核模块} & \shortstack[l]{\legacyold{用户审核}\\\legacyold{FR-YHSH-0002}} & \legacyold{用户可对他人提交的人工审核申请进行审查，支持参考 AI 检测结果并提交个人审核意见。} \tabularnewline
\hline
 \legacyold{用户审核模块} & \shortstack[l]{\legacyold{评论点赞}\\\legacyold{FR-YHSH-0003}} & \legacyold{支持所有用户对审核任务进行评论和点赞操作，促进平台内交流与反馈。} \tabularnewline
\hline
 \legacyold{用户审核模块} & \shortstack[l]{\legacyold{用户举报}\\\legacyold{FR-YHSH-0004}} & \legacyold{针对恶意评论、异常审核结论或违规内容，用户可进行举报并提交说明，由管理端处理。} \tabularnewline
\hline
用户审核模块 & \shortstack[l]{\iternew{论文人工审核}\\\iternew{FR-YHSH-0005}} & \iternew{用户可对他人提交的全篇论文检测任务进行人工审核，针对 AI 生成段落、事实性问题等内容给出审查意见。} \tabularnewline
\hline
用户审核模块 & \shortstack[l]{\iternew{Review人工审核}\\\iternew{FR-YHSH-0006}} & \iternew{用户可对同行评审 Review 检测任务进行人工审核，判断其是否存在 AI 生成或模板化倾向。} \tabularnewline
\hline
用户审核模块 & \shortstack[l]{\iternew{人工审核结果汇总}\\\iternew{FR-YHSH-0007}} & \iternew{平台对用户提交的审核意见进行统计与汇总，生成最终人工审查结果，用于辅助判定内容真实性。} \tabularnewline
\hline

 \legacyold{学术内容检测模块} & \shortstack[l]{\legacyold{提取图片}\\\legacyold{FR-TPJC-0001}} & \legacyold{从用户上传的 PDF、压缩包等非图像格式文件中提取图像，为后续图像检测提供输入。} \tabularnewline
\hline
 \legacyold{学术内容检测模块} & \shortstack[l]{\legacyold{学术图像AI鉴伪}\\\legacyold{FR-TPJC-0002}} & \legacyold{对学术图像进行自动化检测，识别图像是否存在 AI 生成或造假行为。} \tabularnewline
\hline
 \legacyold{学术内容检测模块} & \shortstack[l]{\legacyold{图像可疑区域细粒度标注}\\\legacyold{FR-TPJC-0003}} & \legacyold{在学术图像检测过程中，对 AI 可疑位置进行细粒度标注与可视化展示。} \tabularnewline
\hline
学术内容检测模块 & \shortstack[l]{\iternew{全篇论文AIGC检测}\\\iternew{FR-LWJC-0001}} & \iternew{支持对全文论文进行 AIGC 检测；用户端展示全文级结论与摘要，模型推断由后端完成。} \tabularnewline
\hline
学术内容检测模块 & \shortstack[l]{\iternew{AI生成段落识别}\\\iternew{FR-LWJC-0002}} & \iternew{标识疑似 AI 生成段落及位置；用户端负责高亮与导航展示，识别逻辑由后端完成。} \tabularnewline
\hline
学术内容检测模块 & \shortstack[l]{\iternew{段落级AI生成概率分析}\\\iternew{FR-LWJC-0003}} & \iternew{输出段落级 AI 生成概率或分档；用户端以表格或可视化呈现，计算由后端完成。} \tabularnewline
\hline
学术内容检测模块 & \shortstack[l]{\iternew{事实性鉴伪子结论生成}\\\iternew{FR-LWJC-0004}} & \iternew{生成事实性鉴伪子结论条目；用户端分组列表展示，推理与规则由后端完成。} \tabularnewline
\hline
学术内容检测模块 & \shortstack[l]{\iternew{事实性鉴伪原因说明}\\\iternew{FR-LWJC-0005}} & \iternew{为子结论提供原因与依据说明；用户端展示与折叠交互，文案与证据结构由后端下发。} \tabularnewline
\hline
学术内容检测模块 & \shortstack[l]{\iternew{同行评审Review检测}\\\iternew{FR-PLJC-0001}} & \iternew{对同行评审 Review 内容进行检测，判断其是否存在 AI 生成或异常生成倾向。} \tabularnewline
\hline
学术内容检测模块 & \shortstack[l]{\iternew{模板化评审倾向识别}\\\iternew{FR-PLJC-0002}} & \iternew{针对同行评审 Review 中的模板化表述进行识别，用于发现虚假评审问题。} \tabularnewline
\hline
学术内容检测模块 & \shortstack[l]{\iternew{综合鉴伪结果生成}\\\iternew{FR-NRJC-0001}} & \iternew{系统整合图像、论文与 Review 的检测结果，生成统一的综合鉴伪结论。} \tabularnewline
\hline
学术内容检测模块 & \shortstack[l]{\iternew{检测结果可视化解释}\\\iternew{FR-NRJC-0002}} & \iternew{通过图表、高亮等方式解释检测依据；用户端负责可视化渲染与交互，数据与结论字段由后端提供。} \tabularnewline
\hline
学术内容检测模块 & \shortstack[l]{\iternew{多模态联合检测}\\\iternew{FR-NRJC-0003}} & \iternew{在原有检测能力基础上，新增对论文文本、学术图像及同行评审 Review 的统一联合检测，支持跨模态一致性分析与风险识别。} \tabularnewline
\hline
学术内容检测模块 & \shortstack[l]{\iternew{综合结果融合分析}\\\iternew{FR-NRJC-0004}} & \iternew{对文本、图像及 Review 检测结果进行统一融合分析，输出综合可信度评分、风险等级及最终鉴伪结论，支持生成统一检测报告。} \tabularnewline
\hline

 \legacyold{管理维护模块} & \shortstack[l]{\legacyold{实体添加}\\\legacyold{FR-GLWH-0001}} & \legacyold{在管理端界面可以直接添加图像、论文、Review、用户等实体，新增实体可在用户端显示或使用。} \tabularnewline
\hline
 \legacyold{管理维护模块} & \shortstack[l]{\legacyold{实体修改}\\\legacyold{FR-GLWH-0002}} & \legacyold{在管理端界面可以直接修改图像、论文、Review、用户等实体的信息，修改结果同步至用户端。} \tabularnewline
\hline
 \legacyold{管理维护模块} & \shortstack[l]{\legacyold{实体删除}\\\legacyold{FR-GLWH-0003}} & \legacyold{在管理端界面可以选取指定的图像、论文、Review、用户等实体删除，删除后实体不再显示且不可继续使用。} \tabularnewline
\hline
 \legacyold{管理维护模块} & \shortstack[l]{\legacyold{实体检索}\\\legacyold{FR-GLWH-0004}} & \legacyold{在管理端界面可以对图像、论文、Review、用户等实体进行检索，并支持按标签或类型分类检索。} \tabularnewline
\hline
 \legacyold{管理维护模块} & \shortstack[l]{\legacyold{实体分类}\\\legacyold{FR-GLWH-0005}} & \legacyold{在管理端界面可以对学术资源按照学科、资源类型、造假程度等标签进行分类管理。} \tabularnewline
\hline
 \legacyold{管理维护模块} & \shortstack[l]{\legacyold{数据信息统计}\\\legacyold{FR-GLWH-0006}} & \legacyold{管理端提供多维度数据统计功能，以便管理员了解平台运行情况，包括用户、图像、论文、Review、检测任务与审核任务等数据。} \tabularnewline
\hline
 \legacyold{管理维护模块} & \shortstack[l]{\legacyold{操作记录统计}\\\legacyold{FR-GLWH-0007}} & \legacyold{管理端提供详细的系统日志记录，包括用户操作记录、检测任务记录与审核任务记录等，用于问题排查与系统优化。} \tabularnewline
\hline
 \legacyold{管理维护模块} & \shortstack[l]{\legacyold{超级管理员}\\\legacyold{FR-GLWH-0008}} & \legacyold{超级管理员除了拥有一般管理员的权限外，还可以增加、删除一般管理员并进行权限管理。} \tabularnewline
\hline
 \legacyold{管理维护模块} & \shortstack[l]{\legacyold{用户举报审核}\\\legacyold{FR-GLWH-0009}} & \legacyold{处理用户的举报信息，给予举报方及被举报方做出相应反馈，惩罚措施可包括警告、限制使用、封号等操作。} \tabularnewline
\hline
管理维护模块 & \shortstack[l]{学术资源管理\\FR-GLWH-0010} & 管理员可对平台上传的学术资源（论文、图像、Review）进行全面管理，包括分类、存储、删除等操作。 \tabularnewline
\hline
管理维护模块 & \shortstack[l]{\iternew{论文检测日志管理}\\\iternew{FR-GLWH-0011}} & \iternew{管理端支持记录与管理论文检测相关日志，便于追踪任务执行情况和排查异常。} \tabularnewline
\hline
管理维护模块 & \shortstack[l]{\iternew{Review检测日志管理}\\\iternew{FR-GLWH-0012}} & \iternew{管理端支持记录与管理同行评审 Review 检测相关日志，便于追踪任务执行情况和排查异常。} \tabularnewline
\hline
管理维护模块 & \shortstack[l]{\iternew{模型管理与配置}\\\iternew{FR-GLWH-0013}} & \iternew{后台支持不同检测大模型的配置和管理，便于平台根据任务需求调用不同模型。} \tabularnewline
\hline

\end{longtable}
\end{center}